\section{Взаимодействия и потоки выполнения}
\label{sec:interactions}

В этом разделе фиксируются ключевые сценарии выполнения в вычислительной модели: поиск и применение метода вычисления характеристики, генерация выборки, а также точки расширения (регистрация новых правил и подмена стратегий).

\subsection{Поток вычисления характеристики}

Запрос характеристики начинается с обращения к объекту распределения через \texttt{query\_method} или \texttt{calculate\_characteristic}. Дальнейшее поведение делегируется стратегии вычислений:\begin{enumerate}[compact]
      \item \textbf{Аналитический якорь.} Если в \texttt{analytical\_computations} присутствует реализация для целевой характеристики, стратегия возвращает \texttt{AnalyticalComputation}.
      \item \textbf{Кэш подогнанных методов.} При включённом кэшировании выполняется поиск уже построенного \texttt{FittedComputationMethod} для данной характеристики (и, при необходимости, набора \texttt{**options}). Если метод найден, он возвращается немедленно.
      \item \textbf{Построение по реестру.} Если аналитической реализации нет, стратегия запрашивает у глобального \texttt{CharacteristicRegistry} локальный вид графа: \texttt{view = registry.view(distr)}. В \texttt{RegistryView} выполняется поиск пути от доступных якорей (как правило, от характеристик, присутствующих в \texttt{analytical\_computations}) к целевой характеристике. Поиск пути (\texttt{find\_path}) выполняется с учётом:
            \begin{itemize}[compact]
                  \item наличия узлов и рёбер в локальном виде графа;
                  \item ограничений применимости (constraints), зависящих от признаков типа распределения и/или параметров экземпляра;
                  \item предпочтений по метке варианта (\texttt{prefer\_label}) и глобального порядка \texttt{label\_preference}.
            \end{itemize}
      \item \textbf{Подгонка и композиция.} Для каждого шага найденного пути выбирается конкретный вариант перехода (\texttt{EdgeMeta}), после чего соответствующий \texttt{ComputationMethod} подгоняется под распределение через \texttt{fit(distr, **options)}. Полученная последовательность \texttt{FittedComputationMethod} компонуется в единый вызываемый объект для целевой характеристики.
      \item \textbf{Мемоизация.} Итоговый подогнанный метод сохраняется в кэш стратегии и возвращается вызывающей стороне.
\end{enumerate}

После получения метода \texttt{calculate\_characteristic} применяет его к входному аргументу (например, к точке $x$ для \texttt{pdf}/\texttt{cdf} или к массиву $p$ для \texttt{ppf}).

\subsection{Опции вычислений (\texorpdfstring{\texttt{**options}}{**options})}

Параметр \texttt{**options} используется как сквозной канал настройки вычислений и может влиять на:
\begin{itemize}[compact]
      \item выбор варианта перехода в реестре (например, предпочтение численного/аналитического метода);
      \item поведение подгонки (\texttt{fit}) и подготовку внутренних структур;
      \item поведение вызова подогнанного метода (например, точность, ограничения итераций и т.п.).
\end{itemize}

Архитектурно \texttt{**options} трактуется как часть контекста вычисления: при наличии кэширования рекомендуется включать его (или его нормализованное представление) в ключ кэша.

\subsection{Поток генерации выборки}

Генерация выборки делегируется стратегии сэмплирования:\begin{enumerate}[compact]
      \item вызывается \texttt{distr.sample(n, **options)};
      \item внутри \texttt{sample} выполняется \texttt{sampling\_strategy.sample(n, distr, **options)};
      \item результат возвращается в виде объекта, реализующего протокол \texttt{Sample} (обычно \texttt{ArraySample}).
\end{enumerate}

Такое разделение позволяет:
\begin{itemize}[compact]
      \item подключать специальные алгоритмы сэмплирования для отдельных распределений;
      \item переиспользовать представление выборки и последующую обработку в прикладных пакетах;
      \item тестировать стратегии независимо от конкретных семейств.
\end{itemize}

\subsection{Точки расширения}

Архитектура допускает расширение без модификации существующих интерфейсов:
\begin{itemize}[compact]
      \item \textbf{Новые правила вывода.} Пользователь может добавить в \texttt{CharacteristicRegistry} новую характеристику или новое правило перехода (\texttt{add\_characteristic}, \texttt{add\_computation}), снабдив его ограничениями применимости.
      \item \textbf{Новые семейства и параметризации.} Через \texttt{ParametricFamilyRegister} (см.~рис.~\ref{fig:uml-families}) регистрируются новые семейства; в рамках семейства добавляются параметризации и ограничения.
      \item \textbf{Подмена стратегий.} Для распределения можно заменить \texttt{ComputationStrategy}/\texttt{SamplingStrategy}, что меняет политику выбора методов, кэширование и способ генерации выборки.
      \item \textbf{Варианты методов.} Для одной пары характеристик могут сосуществовать альтернативные реализации, различающиеся меткой и условиями применимости; стратегия выбирает вариант, используя предпочтения и \texttt{**options}.
\end{itemize}
